\documentclass[11pt]{article}

\usepackage[margin=1in]{geometry}
\usepackage{amsmath}
\usepackage{amssymb}
\usepackage{booktabs}
\usepackage{enumitem}
\usepackage{xcolor}
\usepackage[colorlinks=true,linkcolor=blue,citecolor=blue,urlcolor=blue]{hyperref}

\setlength{\parskip}{0.5em}
\setlength{\parindent}{0pt}

\newcommand{\code}[1]{\texttt{#1}}
\newcommand{\rel}[1]{\textsf{#1}}

\title{\textbf{Relational-Geometric Learning over the\\Software Development Lifecycle}\\[0.4em]
\large A Technical Report on a Public Research Lab for\\Verifiable, Relation-Aware SDLC Retrieval}
\author{Relational SDLC AI Lab}
\date{Formalization report --- compiled from the public repository}

\begin{document}
\maketitle

\begin{abstract}
This report formalizes the concepts, data pipeline, models, experimental
protocol, and results of a public research lab studying \emph{relational-geometric
learning} over software-development-lifecycle (SDLC) records. The thesis is that
software-engineering systems should learn and operate over \emph{verifiable
relations} among artifacts (issue, pull request, commit, diff, file, test, log),
not only over the surface similarity of text chunks. We instantiate a typed
artifact/relation graph, a family of relation operators (diagonal metric,
bilinear, low-rank two-tower projection, identity-initialized operator, and a
LoRA-fine-tuned encoder), and a frozen, leakage-guarded benchmark with
Recall@$K$, MRR, and hard-negative metrics. We then report an honest, one-change-
at-a-time experimental arc. On a controlled synthetic benchmark, relation
supervision decisively beats unsupervised IDF and vanilla cosine (Recall@1
\textbf{0.82} vs.\ 0.38 vs.\ 0.11). That win does \emph{not} transfer to real
GitHub issue$\to$PR retrieval, where the surface is rich and IDF is the strong
baseline (R@1 0.46). With references removed and repositories held out,
bag-of-tokens projections fail cross-repo (two-tower 0.24 $<$ vanilla 0.39), while
a frozen pretrained embedder wins outright (0.59), and a bolt-on relation head on
\emph{frozen} vectors cannot improve it. The first positive relational result on
real held-out data is a LoRA fine-tune that places the relation loss \emph{inside}
the representation: R@1 $0.59 \to \textbf{0.66}$, MRR $0.73 \to 0.79$, on eight
unseen repositories, with a $0.48\%$-parameter adapter trained on CPU. The settled
conclusions are that pretrained embeddings are the cross-repo substrate and that
the relational contribution must live in the representation (or in graph
structure), not as a post-hoc head. We close with the next steps. (Several have
since landed --- see the editor's note and the consolidated timeline: the win holds
with multi-split confidence and within-split CIs that exclude zero, grows with
density to $\Delta$R@1 $+0.114$ on a dense 78-repo set, the base axis is
``embedding-tuned'' not ``code,'' and a parameter-free graph lift stacks on top; the
frontier is now GPU/torch-gated scale and a relational small language model.)
\end{abstract}

\tableofcontents
\vspace{1em}

% =====================================================================
\section{Introduction and Thesis}
\label{sec:intro}

Modern software-engineering assistants are built largely on \emph{text} retrieval:
an embedding model maps an artifact to a vector, and relevance is cosine
similarity. This works when the right answer restates the question, but software
artifacts are bound together by \emph{relations} that are often not surface-visible:
a pull request \rel{fixes} an issue, a commit \rel{modifies} a file, a test
\rel{covers} a symbol, a failing log is \rel{caused\_by} a change. These relations
are frequently \emph{verifiable} --- they can be mined from closing keywords, diff
contents, or coverage maps --- which makes them an unusually trustworthy source of
supervision.

\paragraph{Thesis.}
\begin{quote}
Embeddings (and later small language models) trained over \emph{verifiable SDLC
relations} produce more reliable software-engineering systems than models that rely
on generic text similarity.
\end{quote}

The lab's first target is deliberately narrow: not a full autonomous coding agent,
but a \emph{measurable relational retrieval layer}. The target tasks are
\code{issue}$\to$likely fixing PRs; \code{diff}$\to$affected tests; failing
\code{log}$\to$likely files/functions; \code{PR}$\to$risk and missing-test
candidates; and task state$\to$tools or agents to run next. The deliverable is not a
single model but a \emph{frozen public benchmark with honest baselines} that makes
exactly these questions falsifiable.

\paragraph{Five north stars.}
The lab is organized around five standing goals:
(G1) public research hygiene --- public sources and reproducible records;
(G2) falsifiable retrieval gains --- every claim needs a benchmark;
(G3) dataset provenance --- every record carries source, license, timestamp, and
transformation lineage;
(G4) relation-aware models --- models should learn explicit lifecycle relations,
not only text similarity;
(G5) agentic readiness --- every workflow should eventually expose state,
preconditions, actions, tools, and outcomes.

\paragraph{Scope of this report.}
This is a standalone synthesis of the public repository's work. It is
\emph{exploratory} and \emph{pilot-scale} throughout: the numbers are signals to
confirm at scale, not settled results. Where this report and the repository's
experiment cards disagree, the cards are authoritative.

\paragraph{Editor's note (2026-06-23).}
The detailed sections of this \LaTeX{} report (\S\ref{sec:concepts}--\S\ref{sec:results})
were written through the \textbf{Track-A LoRA win} and remain accurate for what they
cover. Subsequent waves --- multi-split confidence (R11A), the graph probe and its
robustness sweep (R11B/R16E), repo-scale and dense Tier-2 (R12A/R13A/R16B/R16C), the
base-model and long-text/chunking studies (R12B/R13B/R14/R15/R16A), the LoRA sweep
(R16D), the headline CIs (R17a), and the diff$\to$test density probe (R17b) --- are
folded into the \emph{consolidated timeline} (\S\ref{sec:results}) and
\emph{What Comes Next} (\S\ref{sec:next}) below, and are covered in full in
\texttt{docs/research-roadmap.md} (\S3 ledger) and the current rendered program report
\texttt{docs/report/relational-sdlc-lab-report-2026-06-23.pdf}. That dated PDF is the
authoritative rendered artifact; this \texttt{.tex} is the \LaTeX{} source of record and
should be recompiled where a \LaTeX{} toolchain is available.

% =====================================================================
\section{Concepts and Formal Framework}
\label{sec:concepts}

\subsection{Typed artifact/relation graph}

The substrate is a typed, directed, attributed multigraph $G = (V, E)$. Nodes are
SDLC \emph{records} drawn from a closed type vocabulary: issues, pull requests,
commits, diffs, files, symbols, tests, CI logs, tools, and agent run logs. Edges
are typed \emph{relations} from the closed set
\[
\mathcal{R} = \{\,\rel{fixes},\ \rel{modifies},\ \rel{covers},\ \rel{fails\_on},\
\rel{caused\_by},\ \rel{requires\_tool},\ \rel{safe\_to\_parallelize},\
\rel{suitable\_agent}\,\}.
\]
Each node carries an embedding-bearing text content plus provenance; each edge
carries an extraction \code{method}, a \code{confidence} $\in [0,1]$, an
\code{observed} flag, and a \code{valid\_from} timestamp used for temporal leakage
control (\S\ref{sec:protocol}). An encoder maps a record to a vector,
\[
E_\theta : \text{record} \longmapsto h \in \mathbb{R}^d,
\]
and a per-relation score function $f_r$ judges whether a directed pair stands in
relation $r$:
\[
s_r(u,v) = f_r(h_u, h_v).
\]
Retrieval for a task is then: given a query record and a fixed candidate pool, rank
candidates by $s_r$.

\subsection{Relation operators}
\label{sec:operators}

The lab studies a ladder of increasingly expressive operators $f_r$, each testing a
different hypothesis about \emph{where} the relational signal lives.

\begin{description}[leftmargin=1.4em,style=nextline]
\item[Translational operator $h_v \approx T_r(h_u)$.]
A relation is a transformation of the source embedding; in the linear case
$T_r(h) = M_r h$, recovering the translational/relation-map family. Implemented as
the identity-initialized operator $M$ (\S\ref{sec:models}).
\item[Bilinear score $s_r(u,v) = h_u^{\top} W_r\, h_v$.]
A learned relation-specific bilinear form. Two restrictions are studied: a
\emph{diagonal} $W_r$ (a per-token reweighting --- it can only emphasize or
suppress \emph{shared} dimensions), and a \emph{low-rank} $W_r = W_q^{\top} W_d$,
which factors into a two-tower projection (it can align \emph{different} dimensions
on the two sides).
\item[Two-tower / low-rank projection.]
Separate projections for the query side and document side,
$a = W_q u_{\text{issue}}$, $b = W_d u_{\text{doc}}$, $\text{score} = \langle a, b
\rangle$. This is the operator capable of \emph{cross-token} structure --- mapping
issue-symptom vocabulary to change vocabulary --- which a diagonal metric cannot
represent.
\end{description}

\subsection{Losses}
\label{sec:losses}

The operators are trained with a \emph{topological margin loss} $L_{\text{topo}}$:
a triplet/contrastive objective that pulls a query toward its related artifact and
pushes it away from a near-miss negative, by at least a margin $\gamma$:
\[
L_{\text{topo}} = \sum_{(u, v^{+}, v^{-})}
\big[\, \gamma - s_r(u, v^{+}) + s_r(u, v^{-}) \,\big]_{+}.
\]
At scale this is realized as a symmetric InfoNCE / multiple-negatives contrastive
loss with in-batch negatives. The full paper program also specifies a supervised
contrastive term $L_{\text{contrastive}}$, a relation-classification term
$L_{\text{rel}}$, a graph-smoothness term $L_{\text{graph}}$, a logical-consistency
term $L_{\text{logic}}$, and an alignment term $L_{\text{align}}$; these remain
partial or future and are noted in \S\ref{sec:next}.

\subsection{Metrics}
\label{sec:metrics}

Every system returns a ranked candidate list per query; aggregated over a frozen
evaluation set we report:
\begin{itemize}[leftmargin=1.4em,nosep]
\item \textbf{Recall@$K$} --- fraction of relevant items in the top $K$
($K \in \{1,5,10\}$);
\item \textbf{MRR} --- mean reciprocal rank of the first relevant item;
\item \textbf{Hard-negative accuracy} --- fraction of queries where a relevant item
outranks every \emph{designated} hard negative (the wrong file in the same package,
the wrong test in the same suite, a plausible-but-unrelated PR).
\end{itemize}
Hard negatives are the point: generic similarity already finds easy positives, so
without near-miss negatives Recall@$K$ flatters every model. Generalization is
measured by \emph{cross-repo} (held-out repositories) and, where relevant,
\emph{temporal} (train on the past, evaluate on the future) splits.

% =====================================================================
\section{Data Pipeline}
\label{sec:data}

The pipeline turns public GitHub data into a provenance-bearing, leakage-guarded,
frozen benchmark. It is methodology-as-code: every record carries provenance, every
edge an extraction method and confidence, every claim an experiment card, every
benchmark a frozen split with a leakage guard.

\subsection{Ingest with provenance}

Ingest is pure standard library and \emph{offline by default}: the default path
replays a recorded API snapshot and never opens a socket; a live fetch is opt-in
behind an explicit flag and a token, is polite by construction (descriptive
user-agent, pinned API version, rate-limit-header obedience, modest pagination),
and is recorded once so the mapping is replayable offline. The mapping is:
\begin{center}
\begin{tabular}{lll}
\toprule
GitHub source & Record / edge & Method \\
\midrule
issue / PR / commit JSON & \code{issue} / \code{pull\_request} / \code{commit} & provenance only \\
PR's changed file & \code{file} or \code{test} record & --- \\
``Fixes/Closes \#N'' in PR body & \rel{fixes} edge & \code{human\_label} ($c{=}0.95$) \\
PR's changed files & \rel{modifies} edge & \code{git\_history} ($c{=}1.0$) \\
\bottomrule
\end{tabular}
\end{center}
Every record and edge carries \code{source\_url}, \code{retrieved\_at},
\code{license}, a real \code{sha256:} content hash, a \code{transform}, a
\code{method}, and an \code{observed} flag. A closing-keyword reference to an issue
\emph{not} in the snapshot is never minted as an edge, so referential integrity
holds. Test files (paths under \code{tests/}, \code{*\_test.*}, \code{*.spec.*},
\dots) are typed as \code{test} records. Redistribution defaults to
\code{metadata\_only}: provenance points back to public URLs; full source text is
not redistributed.

\subsection{Records, edges, and the leakage guard}

The validator enforces JSON schemas, provenance completeness, referential integrity
(every edge endpoint resolves to a record), and a \emph{temporal leakage guard}:
each record/edge carries \code{valid\_from}, and a query's optional \code{as\_of}
time excludes any candidate or positive that only becomes valid later. Random splits
leak; the default split for real data is \emph{temporal-by-commit-date}.

\subsection{De-referenced, frozen, cross-repo benchmarks}

Two design choices make the benchmark honest (rationale in \S\ref{sec:design}):
\textbf{reference removal} --- explicit links (\code{\#123}, \code{owner/repo\#N},
issue/PR URLs, commit SHAs) are scrubbed from \emph{inputs} while kept as ground-
truth labels --- and \textbf{cross-repo splits} --- train repositories are disjoint
from test repositories. Splits are frozen and recorded in a dataset card. Heavy
steps (embedding) are cached and committed so downstream runs replay on the cache
with numpy alone --- no GPU, no network, CI-reproducible.

\subsection{Graph enrichment}

The first pilot snapshot carried \code{issue}/\code{pull\_request}/\code{commit}
records and \rel{fixes} edges but no file-level structure. A one-time live
enrichment fetches each distinct fixing-PR's changed files and maps them into
\code{file}/\code{test} records and \rel{modifies} edges using the same
provenance-bearing helpers, yielding the structure the graph-learning track and the
\code{diff}$\to$affected-test task need. The committed snapshot adds 497 file
records, 239 test records (736 total), and 1{,}356 \rel{modifies} edges over 18
repositories, each edge \code{observed} from the diff with \code{confidence} $1.0$
and temporally consistent endpoints.

% =====================================================================
\section{Models Implemented}
\label{sec:models}

Six systems are scored on the \emph{same} candidate pools so every comparison is
attributable to one change. Each tests a specific hypothesis about where the
relational signal lives.

\begin{description}[leftmargin=1.4em,style=nextline]
\item[Vanilla cosine.] Plain cosine over bag-of-token vectors --- the
off-the-shelf text-similarity floor. With all weights equal to one, every weighted
model below reduces to this, so any gain is purely from supervision.
\item[Unsupervised IDF cosine.] Cosine with inverse-document-frequency token
weights from corpus statistics only --- no relation labels. The bag-of-tokens bar:
if a relation-aware model only matches this, the ``gain'' was just a frequency
effect.
\item[Diagonal relation metric.] Learns a non-negative weight per token via a
margin triplet loss over \emph{train-split} \rel{fixes} pairs. It can up-weight
predictive tokens and down-weight ambiguous ones, but only over \emph{shared}
dimensions. Tests: \emph{is the relation a per-token reweighting?}
\item[Two-tower projection.] Asymmetric low-rank query/document projections trained
with a margin triplet loss. It can align \emph{different} tokens across sides.
Tests: \emph{can a learned cross-token operator (over bag-of-tokens features)
generalize?}
\item[Identity-init operator.] A relation map $M$ \emph{initialized at the
identity} (so it starts exactly at embedder cosine) and regularized back toward it,
applied on \emph{frozen} pretrained embeddings. Tests: \emph{can a safe refinement
of frozen vectors help?}
\item[LoRA-fine-tuned encoder.] Low-rank adapters ($r{=}8$, on attention
query/key/value; $0.48\%$ of parameters) on a frozen small pretrained embedder,
trained with the InfoNCE relation loss on train-repo pairs. The pretrained weights
are untouched; only the adapters learn. Tests: \emph{does the relation loss inside
the representation generalize cross-repo?}
\end{description}

The synthetic generator and the bag-of-tokens models are numpy-only and
deterministic given a seed; the pretrained-embedder systems use a frozen
\code{MiniLM-L6} (22M params, 384-d, CPU), embedded once and cached so the
evaluation is numpy-only.

% =====================================================================
\section{Experimental Protocol and Design Choices}
\label{sec:protocol}

\subsection{Method invariants}

Every experiment obeys six invariants so results are trustworthy and comparable:
\begin{enumerate}[leftmargin=1.6em,nosep]
\item \textbf{Links are labels, not features.} A relation recoverable by regex is a
label source, never a test --- explicit references are scrubbed from inputs.
\item \textbf{Frozen splits; cross-repo is the headline.} Held-out repositories
measure generalization, not repo-specific memorization.
\item \textbf{Pretrained embeddings are the substrate; the relational contribution
is the delta over \code{embedder-cosine}.} A method ``wins'' only if it beats that
control on the same split.
\item \textbf{Every run emits an experiment card}; exploratory results are labeled
exploratory and never presented as release evidence.
\item \textbf{Reproducible from a clean checkout.} Heavy steps are cached and
committed; downstream runs use numpy alone.
\item \textbf{Honest negatives are kept.} A result that refutes a hypothesis is a
result.
\end{enumerate}

\subsection{Design choices, with rationale}
\label{sec:design}

\paragraph{Why links-as-labels (de-reference the input).}
A \rel{fixes} edge stated in the text (``Fixes \#123'') is trivially recoverable by
a regex, so a model that ``solves'' the task by string-matching an issue number has
learned nothing. Scrubbing the reference from the input --- while keeping it as the
ground-truth label --- forces the model to match on \emph{semantics}, making the
benchmark a genuine test of the relational hypothesis rather than a pointer-chasing
exercise.

\paragraph{Why cross-repo.}
The goal is a relation that \emph{generalizes}, not one that memorizes a
repository's idioms. Holding out entire repositories (train repos disjoint from test
repos) is the strict test: a win must survive unseen vocabularies, module names, and
API surfaces. This is exactly the split on which bag-of-tokens projections were
found to fail.

\paragraph{Why pretrained embeddings as the substrate.}
The cross-repo experiments showed, with evidence, that \emph{tokens do not transfer
between repositories; only meaning does}. A from-scratch token projection assigns
random columns to exactly the discriminative vocabulary it never saw in training. A
pretrained semantic embedder already places related meanings near each other across
repositories, so it is the right feature substrate on which to learn relations.

\paragraph{Why LoRA (and why the from-scratch head failed).}
On \emph{frozen} embeddings at pilot scale, a bolt-on relation head has no headroom:
a from-scratch projection overfits the train repos and \emph{destroys} the
embedder's native geometry, while an identity-init operator is safe but adds
nothing. The conclusion is that the relational contribution must \emph{reshape} the
representation rather than project a frozen one. LoRA is the disciplined way to do
that: tiny, low-rank adapters move the representation with very few parameters and
little compute, keeping the pretrained weights intact.

\paragraph{The role of honest negatives.}
Keeping refuting results (the non-transfer to real data; the cross-repo failure of
bag-of-tokens; the no-headroom of frozen-embedder heads) is not modesty --- it is
what \emph{narrowed the search}. Each negative eliminated a hypothesis and pointed
at the next experiment, culminating in the place the contribution had to live.

\subsection{Compute discipline}

CPU-first; small models (a small embedder now, a small SLM later); LoRA, not full
fine-tuning; cache and commit embeddings so experiments replay cheaply; reach for
GPU/scale only when a method has shown signal worth scaling. The benchmark stays
small until the signal says ``feed me.''

% =====================================================================
\section{Results}
\label{sec:results}

The work proceeded as a sequence of one-change-at-a-time experiments, each against
the appropriate control. Numbers are pulled from the ablation documents and the
committed experiment cards; all are \emph{exploratory, pilot-scale}.

\subsection{Synthetic per-token ablation (mechanism)}

A controlled benchmark where the \rel{fixes} link is carried by \emph{rare,
component-specific} ``impl'' tokens while the surface is dominated by \emph{common,
ambiguous} ``topic'' tokens. Hard negatives are PRs from other components that share
topic words. On 94 held-out queries:
\begin{center}
\begin{tabular}{lcccc}
\toprule
System & Recall@1 & Recall@5 & MRR & Hard-neg \\
\midrule
vanilla-tf-cosine & 0.106 & 0.309 & 0.276 & 0.106 \\
idf-cosine & 0.383 & 0.532 & 0.501 & 0.383 \\
\textbf{relation-metric} & \textbf{0.819} & \textbf{0.894} & \textbf{0.864} & \textbf{0.819} \\
\bottomrule
\end{tabular}
\end{center}
The learned weights tell the story: mean impl-token weight $\approx 1.36$, mean
topic-token weight $\approx 0.28$ --- the supervision recovered which tokens predict
the fix, without being told. \emph{Finding: confirmed} --- relation supervision
recovers a link that surface similarity cannot, on a controlled benchmark.

\subsection{Cross-token synthetic (operator validation)}

A harder synthetic split where the issue and its fix share \emph{no} tokens. Vanilla,
IDF, and the diagonal metric all sit at chance (R@1 $0.01$ --- they can only reweight
shared tokens, of which there are none); the two-tower operator recovers the mapping
at R@1 \textbf{0.832}. \emph{Finding: confirmed} --- only a cross-token operator can
bridge disjoint vocabulary; the tower implementation is correct.

\subsection{Real issue$\to$PR (the synthetic win does not transfer)}

The P1 pilot --- a frozen snapshot of 20 permissive Python-ecosystem repos
(2{,}087 records, 356 \rel{fixes} edges, 356 queries, temporal split). On 144
held-out queries:
\begin{center}
\begin{tabular}{lcccc}
\toprule
System & Recall@1 & Recall@5 & MRR & Hard-neg \\
\midrule
vanilla-tf-cosine & 0.347 & 0.715 & 0.522 & 0.354 \\
\textbf{idf-cosine} & \textbf{0.438} & \textbf{0.812} & \textbf{0.611} & 0.438 \\
relation-metric & 0.417 & 0.785 & 0.595 & 0.438 \\
\bottomrule
\end{tabular}
\end{center}
\emph{Finding: no transfer.} Real PRs restate the issues they fix, so the surface is
already rich (vanilla R@5 $\approx 0.72$); IDF is the strong baseline and the diagonal
metric ties it. A per-token reweighting has nothing extra to learn --- real
issue$\leftrightarrow$fix association is cross-token and semantic.

\subsection{De-referenced, cross-repo, bag-of-tokens}

References removed and repositories held out (10 train / 8 test). On 174 held-out-repo
queries:
\begin{center}
\begin{tabular}{lcccc}
\toprule
System & Recall@1 & Recall@5 & MRR & Hard-neg \\
\midrule
vanilla-tf-cosine & 0.391 & 0.724 & 0.544 & 0.397 \\
\textbf{idf-cosine} & \textbf{0.460} & \textbf{0.828} & \textbf{0.624} & 0.460 \\
relation-metric (diagonal) & 0.454 & 0.799 & 0.618 & 0.471 \\
relation-tower (projection) & 0.241 & 0.713 & 0.445 & 0.362 \\
\bottomrule
\end{tabular}
\end{center}
\emph{Finding: no.} The tower --- validated on cross-token synthetic --- scores
\emph{below vanilla} (0.241) cross-repo, because the discriminative vocabulary of
unseen repos was never trained, so its projection columns are random on exactly the
tokens that matter. IDF stays the robust winner precisely because it does not train.
\emph{Tokens do not transfer between repositories; only meaning does.}

\subsection{Pretrained embeddings settle the substrate}

The same de-referenced cross-repo split, now with a frozen \code{MiniLM-L6} embedder
(zero training):
\begin{center}
\begin{tabular}{lcccc}
\toprule
System & Recall@1 & Recall@5 & MRR & Hard-neg \\
\midrule
idf-cosine (bag-of-tokens bar) & 0.460 & 0.828 & 0.624 & 0.460 \\
\textbf{embedder-cosine} & \textbf{0.592} & 0.920 & 0.728 & 0.592 \\
embedder + from-scratch head & 0.190 & 0.644 & 0.381 & 0.218 \\
embedder + identity-init operator & 0.592 & \textbf{0.931} & \textbf{0.737} & 0.592 \\
\bottomrule
\end{tabular}
\end{center}
\emph{Findings.} (1) Pretrained embeddings are the cross-repo win (R@1 $0.46 \to
0.59$, no training) --- the substrate question is settled. (2) A from-scratch head on
frozen embeddings is actively \emph{harmful} (0.19): it overfits and destroys the
embedder's geometry. (3) An identity-init operator is safe but adds essentially
nothing at this scale. \emph{The relational contribution cannot be a post-hoc head on
frozen vectors.}

\subsection{The relational fine-tune wins (Track A)}

Putting the contribution \emph{inside} the representation: LoRA adapters ($r{=}8$,
$0.48\%$ of params) on the same embedder, trained with the InfoNCE relation loss on
182 train-repo pairs, evaluated on the same eight held-out repos:
\begin{center}
\begin{tabular}{lcccc}
\toprule
System & Recall@1 & Recall@5 & MRR & Hard-neg \\
\midrule
idf-cosine (bag-of-tokens bar) & 0.460 & 0.828 & 0.624 & 0.460 \\
embedder-cosine (frozen) & 0.592 & 0.920 & 0.728 & 0.592 \\
\textbf{embedder-cosine (LoRA-tuned)} & \textbf{0.655} & \textbf{0.994} & \textbf{0.791} & \textbf{0.661} \\
LoRA + identity-init operator & 0.649 & 0.994 & 0.790 & 0.655 \\
\bottomrule
\end{tabular}
\end{center}
\emph{Findings.} (1) The relational fine-tune \emph{wins cross-repo}: R@1 $0.592 \to
\mathbf{0.655}$, MRR $0.728 \to \mathbf{0.791}$, on repositories not seen in training
--- the first positive relational result on real held-out data. (2) The contribution
is in the representation, not a head: adding the identity-init operator on top of the
\emph{tuned} vectors changes nothing ($0.649 \approx 0.655$), exactly as predicted.
(3) A $0.48\%$-param adapter on a 22M model, trained on 182 pairs on CPU, moved the
needle --- evidence the signal is real and learnable, not a capacity artifact.

\subsection{Consolidated timeline}

\begin{center}
\begin{tabular}{p{4.6cm}p{4.4cm}cc}
\toprule
Experiment & Finding & R@1 (winner) & vs.\ control \\
\midrule
Synthetic per-token & relation supervision wins & \textbf{0.82} & vs.\ IDF 0.38 / van.\ 0.11 \\
Cross-token synthetic & only the tower bridges disjoint vocab & \textbf{0.83} & vs.\ all-else 0.01 \\
Real issue$\to$PR (linked) & no transfer; surface-rich & IDF \textbf{0.46} & diag.\ ties (0.45) \\
De-ref.\ cross-repo (tokens) & bag-of-tokens fails cross-repo & tower \textbf{0.24} & $<$ vanilla 0.39 \\
Embeddings cross-repo & pretrained embeddings win & \textbf{0.59} & vs.\ IDF 0.46 \\
Head on frozen embeddings & no headroom; head harmful & 0.19 / 0.59 & ties / hurts \\
LoRA fine-tune (Track A) & \textbf{relational win in the rep} & \textbf{0.66} & vs.\ frozen 0.59 \\
\midrule
\multicolumn{4}{l}{\emph{\small Waves after the Track-A win (full ledger: \texttt{research-roadmap.md} \S3):}} \\
Multi-split (R11A) & LoRA win on \emph{all 5} held-out splits & $\Delta$\textbf{+0.061} & $\pm0.021$ \\
Graph probe (R11B) & free aggregation stacks on LoRA & \textbf{0.69} & vs.\ LoRA 0.66 \\
Dense Tier-2 (R16C) & LoRA win \emph{grows} with density & $\Delta$\textbf{+0.114} & $0.515{\to}0.629$ \\
Base model (R12B/R13B) & ``embedding-tuned'', not ``code'' & 0.14--0.60 & monotone in tuning \\
Long text / chunk (R14/R15/R16A) & truncation helps issue$\to$PR; MaxP wins deep & --- & FirstP@512 0.701 \\
Graph sweep (R16E) & robust plateau; diff$\to$test structure-bound & 0.690 & ceiling 59.8\% \\
LoRA-win CIs (R17a) & 95\% CIs exclude zero (query + repo) & $\Delta$+0.063 & CI [+.006,+.121] \\
diff$\to$test density (R17b) & 59.8\% ceiling is an ingest artefact & \textbf{96.4\%} & co-change median 35 \\
\bottomrule
\end{tabular}
\end{center}

% =====================================================================
\section{Discussion}
\label{sec:discussion}

\paragraph{What is settled.}
Two conclusions are evidence-backed at pilot scale.
\emph{(i) Pretrained embeddings are the cross-repo substrate.} Meaning transfers
across repositories where tokens do not; a frozen semantic embedder beats every
bag-of-tokens system on held-out repos with no training.
\emph{(ii) The relational contribution belongs in the representation (or in graph
structure), not in a post-hoc head.} A bolt-on operator on frozen vectors has no
headroom and can be harmful, whereas reshaping the encoder with the relation loss
(LoRA) generalizes cross-repo. These two findings jointly redirect the program: the
``relational contribution as the core'' is a fine-tuned encoder, not a reranker on
fixed features.

\paragraph{What is open.}
The Track-A win is a single seed, a single cross-repo split, 182 train pairs, and a
general-text base; R@5 $\approx 0.99$ is near the ceiling of small candidate pools, so
R@1 and MRR are the discriminating metrics. Open questions remain falsifiable:
whether the win survives re-splits and multiple seeds; whether graph link prediction
adds signal beyond pairwise cosine; whether relations with weaker surface text
(\code{diff}$\to$test, \code{log}$\to$file) show a larger lift; whether cross-repo
same-bug-class retrieval benefits; whether scale gives a learned head the headroom it
lacked; and how much a code-specific base helps.

\paragraph{Method as the deliverable.}
The lab's enduring output is not any one number but the \emph{frozen public benchmark
with honest baselines and full provenance} that makes these questions measurable. The
embedding layer is built to earn trust first; the relational small language model and
agent policy come only after retrieval, provenance, and evaluation gates are boring.

% =====================================================================
\section{What Comes Next}
\label{sec:next}

The decision gate after the Track-A win routes to two parallel directions, with later
phases gated on signal.

\begin{description}[leftmargin=1.4em,style=nextline]
\item[Scale the win (Track D).] \emph{Largely done and now GPU-gated.} The LoRA result
is confirmed with multi-split confidence (R11A) and within-split CIs that exclude zero
(R17a); it \emph{grows} with density through a dense 78-repo Tier-2 set
($\Delta$R@1 $+0.114$, R16C); the base axis is ``embedding-tuned,'' not ``code''
(R12B/R13B); and a rank sweep shows the Tier-2 delta is near-tuned, with harder in-batch
negatives the only remaining lever (R16D). What remains needs a GPU: Tier-2 full breadth
(200--500 repos), larger negative pools, and multiple seeds.
\item[Graph link prediction (Track B).] The graph-enrichment prerequisite is done
(file/test nodes and \rel{modifies} edges). The next step is an inductive GNN
(GraphSAGE / R-GCN) on \emph{pretrained} node features so it transfers to unseen
repos, with multi-relation link prediction, plus KG-embedding scoring (DistMult /
ComplEx / RotatE) for asymmetric relations. The control is \code{embedder-cosine}
plus the Track-A fine-tune; a win must beat them on a relation cosine cannot capture.
\item[Tasks and labels (Track C).] Activate \code{diff}$\to$affected-test and
\code{log}$\to$file (Q3) on the file edges, and cross-repo same-bug-class retrieval
(Q4) using a curated SWE bug-class ontology --- with labels independent of the body
text to avoid circularity.
\item[Remaining losses.] Complete $L_{\text{contrastive}}$, $L_{\text{rel}}$,
$L_{\text{graph}}$, $L_{\text{logic}}$, and $L_{\text{align}}$ as the model family
matures.
\item[Relational SLM (MVP-2+, Track E).] Only after MVP-1 holds: a QLoRA small
language model with relation/policy heads for review, test suggestion, and risk; a
GraphRAG subgraph packer; and an agent policy with outcome learning. Autonomous
editing comes last, after retrieval and policy heads are measurable.
\end{description}

% =====================================================================
\section{References}
\label{sec:refs}

\emph{Cited by name; arXiv-style, public methods only.}

\begin{enumerate}[leftmargin=1.6em,nosep]
\item F.\ Schroff, D.\ Kalenichenko, J.\ Philbin. \emph{FaceNet: A Unified Embedding
for Face Recognition and Clustering.} CVPR, 2015. (triplet margin loss)
\item P.\ Khosla et al. \emph{Supervised Contrastive Learning} (SupCon). NeurIPS,
2020.
\item N.\ Reimers, I.\ Gurevych. \emph{Sentence-BERT: Sentence Embeddings using
Siamese BERT-Networks.} EMNLP, 2019.
\item T.\ Gao, X.\ Yao, D.\ Chen. \emph{SimCSE: Simple Contrastive Learning of
Sentence Embeddings.} EMNLP, 2021.
\item W.\ Wang et al. \emph{MiniLM: Deep Self-Attention Distillation for
Task-Agnostic Compression of Pre-Trained Transformers.} NeurIPS, 2020. (the
\code{all-MiniLM-L6-v2} base)
\item B.\ Yang, W.\ Yih, X.\ He, J.\ Gao, L.\ Deng. \emph{Embedding Entities and
Relations for Learning and Inference in Knowledge Bases} (DistMult). ICLR, 2015.
\item T.\ Trouillon et al. \emph{Complex Embeddings for Simple Link Prediction}
(ComplEx). ICML, 2016.
\item Z.\ Sun, Z.\ Deng, J.\ Nie, J.\ Tang. \emph{RotatE: Knowledge Graph Embedding
by Relational Rotation in Complex Space.} ICLR, 2019.
\item T.\ Kipf, M.\ Welling. \emph{Semi-Supervised Classification with Graph
Convolutional Networks} (GCN). ICLR, 2017.
\item W.\ Hamilton, R.\ Ying, J.\ Leskovec. \emph{Inductive Representation Learning
on Large Graphs} (GraphSAGE). NeurIPS, 2017.
\item M.\ Schlichtkrull et al. \emph{Modeling Relational Data with Graph
Convolutional Networks} (R-GCN). ESWC, 2018.
\item E.\ Hu et al. \emph{LoRA: Low-Rank Adaptation of Large Language Models.} ICLR,
2022.
\item T.\ Dettmers, A.\ Pagnoni, A.\ Holtzman, L.\ Zettlemoyer. \emph{QLoRA:
Efficient Finetuning of Quantized LLMs.} NeurIPS, 2023.
\item P.\ Lewis et al. \emph{Retrieval-Augmented Generation for Knowledge-Intensive
NLP Tasks} (RAG). NeurIPS, 2020.
\item B.\ Peng et al. \emph{GraphRAG: Graph Retrieval-Augmented Generation} (survey
of graph-based RAG), 2024.
\item C.\ Jimenez et al. \emph{SWE-bench: Can Language Models Resolve Real-World
GitHub Issues?} ICLR, 2024.
\end{enumerate}

\vfill
\begin{center}
\small\emph{Exploratory, pilot-scale results. The committed experiment cards are the
source of truth for every number in this report.}
\end{center}

\end{document}
